\documentclass[tikz,border=12pt]{standalone}
\usepackage{amsmath}
\usepackage{kotex}
\usetikzlibrary{arrows.meta,calc,positioning}

\definecolor{stress}{RGB}{198,68,54}
\definecolor{strain}{RGB}{40,105,175}
\definecolor{body}{RGB}{225,235,244}
\definecolor{edge}{RGB}{55,70,85}

\begin{document}
\begin{tikzpicture}[>=Latex,font=\sffamily,line cap=round,line join=round]
  \node[font=\Large\bfseries] at (7.2,7.4) {평면응력과 평면변형률: 면외에서 무엇을 0으로 두는가?};

  % 왼쪽: 평면응력
  \node[font=\large\bfseries,stress] at (3.3,6.55) {평면응력};
  \node[font=\small] at (3.3,6.08) {얇은 판의 중앙면};

  \fill[body] (0.8,2.6) -- (5.6,2.6) -- (6.25,3.05) -- (1.45,3.05) -- cycle;
  \fill[body!70!white] (0.8,2.6) -- (1.45,3.05) -- (1.45,3.35) -- (0.8,2.9) -- cycle;
  \fill[body!85!white] (1.45,3.05) -- (6.25,3.05) -- (6.25,3.35) -- (1.45,3.35) -- cycle;
  \draw[edge,thick] (0.8,2.6) -- (5.6,2.6) -- (6.25,3.05) -- (1.45,3.05) -- cycle;
  \draw[edge,thick] (0.8,2.9) -- (1.45,3.35) -- (6.25,3.35) -- (6.25,3.05);
  \draw[edge,thick] (0.8,2.6) -- (0.8,2.9) -- (1.45,3.35);

  \draw[stress,very thick,->] (0.8,2.75) -- (-0.25,2.75);
  \draw[stress,very thick,->] (5.9,3.18) -- (6.95,3.18);
  \draw[stress,very thick,->] (3.55,2.62) -- (3.55,1.65);
  \draw[stress,very thick,->] (3.55,3.35) -- (3.55,4.32);
  \node[stress,font=\small] at (3.55,4.65) {면내 하중};

  \draw[->,thick] (1.25,1.35) -- (2.15,1.35) node[right] {$\boldsymbol e_1$};
  \draw[->,thick] (1.25,1.35) -- (1.75,1.85) node[above] {$\boldsymbol e_2$};
  \draw[->,thick] (1.25,1.35) -- (1.25,2.22) node[above] {$\boldsymbol e_3$};
  \node[align=center,font=\small] at (3.55,0.72)
    {자유표면의 면외 면력 $\approx 0$\\
     \textcolor{stress}{$\sigma_{33}=\sigma_{13}=\sigma_{23}\approx0$}\\
     \textcolor{strain}{$\varepsilon_{33}$은 일반적으로 $0$이 아님}};

  % 구분선
  \draw[gray!45,dashed] (7.2,0.1) -- (7.2,6.65);

  % 오른쪽: 평면변형률
  \node[font=\large\bfseries,strain] at (10.95,6.55) {평면변형률};
  \node[font=\small] at (10.95,6.08) {긴 구조물의 중앙 단면};

  \fill[body] (8.2,2.25) rectangle (13.5,4.15);
  \fill[body!72!white] (13.5,2.25) -- (14.1,2.75) -- (14.1,4.65) -- (13.5,4.15) -- cycle;
  \fill[body!88!white] (8.2,4.15) -- (13.5,4.15) -- (14.1,4.65) -- (8.8,4.65) -- cycle;
  \draw[edge,thick] (8.2,2.25) rectangle (13.5,4.15);
  \draw[edge,thick] (8.2,4.15) -- (8.8,4.65) -- (14.1,4.65) -- (13.5,4.15);
  \draw[edge,thick] (13.5,2.25) -- (14.1,2.75) -- (14.1,4.65);

  \fill[strain!28] (10.6,2.25) rectangle (11.05,4.15);
  \draw[strain,very thick] (10.6,2.25) rectangle (11.05,4.15);
  \node[strain,font=\small,align=center] at (10.83,5.25) {관심 있는\\중앙 단면};
  \draw[strain,->,thick] (10.83,4.95) -- (10.83,4.2);

  \draw[->,thick] (8.55,1.3) -- (9.45,1.3) node[right] {$\boldsymbol e_3$};
  \draw[->,thick] (8.55,1.3) -- (8.55,2.15) node[above] {$\boldsymbol e_2$};
  \draw[->,thick] (8.55,1.3) -- (9.05,1.78) node[above right] {$\boldsymbol e_1$};

  \draw[strain,very thick,->] (8.25,1.72) -- (13.45,1.72);
  \node[strain,font=\small,fill=white,inner sep=1pt] at (10.85,1.72) {$\boldsymbol e_3$ 방향: 변위장이 거의 변하지 않음};
  \node[align=center,font=\small] at (10.95,0.55)
    {면외 변형이 억제됨\\
     \textcolor{strain}{$\varepsilon_{33}=\varepsilon_{13}=\varepsilon_{23}\approx0$}\\
     \textcolor{stress}{$\sigma_{33}$은 일반적으로 $0$이 아님}};

  \node[align=center,text=gray!80!black,font=\small] at (7.2,-0.62)
    {두 조건 모두 3차원 문제의 근사이다. 평면응력은 면외 \textbf{응력}을, 평면변형률은 면외 \textbf{변형률}을 0으로 둔다.};
\end{tikzpicture}
\end{document}
